% 01 - CONTEXTE
% ==============================================================
\sectionintro{01 · Démarche}{Contexte et objectif}{Conserver la trace de la première approche, puis décrire fidèlement l'état actuel du projet TwinCAT.}

\noindent
\begin{minipage}[t]{0.355\textwidth}\vspace{0pt}
  \begin{primarycard}
    \textbf{Positionnement}\par\small
    Lors du test, une version volontairement simplifiée a été développée afin de présenter rapidement le principe de fonctionnement. Le projet a ensuite été repris sous TwinCAT pour approfondir l'architecture, la gestion des huit bacs, les réinitialisations et plusieurs cas limites.
  \end{primarycard}

  \subsection{Objectif du rapport}
  \small
  Le document distingue la \textbf{solution de principe montrée pendant l'entretien} de la \textbf{version approfondie actuellement présente dans le dépôt}. Les résultats d'essais déjà consignés sont conservés comme tels ; la présente mise à jour vérifie statiquement les objets TwinCAT et n'affirme pas avoir rejoué le runtime. Les fonctions non réalisées sont identifiées comme limites ou évolutions.

  \subsection{Fil conducteur}
  \begin{enumerate}\small
    \item poser le cycle nominal ;
    \item isoler les modes et les états ;
    \item traduire la séquence en ST ;
    \item provoquer des cas limites ;
    \item comparer le rapport aux sources PLC et à la Visualization V1.
  \end{enumerate}
\end{minipage}\hfill
\begin{minipage}[t]{0.605\textwidth}\vspace{0pt}
  \begin{card}
    \textbf{Première lecture fonctionnelle}\par\small
    Le schéma de séquencement initial reste une trace utile : il montre le raisonnement présenté à l'automaticien sans présenter cette première version comme aussi complète que le projet final.
  \end{card}

  \begin{center}
  \begin{tikzpicture}[node distance=8mm and 9mm]
    \node[ginit] (s0) {0};
    \node[gstep,right=of s0] (s1) {1};
    \node[gstep,right=of s1] (s2) {2};
    \node[gstep,right=of s2,fill=eksoft] (s3) {3};
    \node[gstep,right=of s3] (s4) {4};
    \draw[flow] (s0)--(s1);
    \draw[flow] (s1)--(s2);
    \draw[flow] (s2)--(s3);
    \draw[flow] (s3)--(s4);
    \draw[flow] (s4.south)--++(0,-7mm)-|(s1.south);
    \node[glabel,above=5mm of s0,text width=18mm,align=center] {initialiser};
    \node[glabel,above=5mm of s1,text width=18mm,align=center] {compter};
    \node[glabel,above=5mm of s2,text width=18mm,align=center] {tester};
    \node[glabel,above=5mm of s3,text width=18mm,align=center] {indexer};
    \node[glabel,above=5mm of s4,text width=18mm,align=center] {bac suivant};
  \end{tikzpicture}
  \end{center}

  \begin{notebox}{Ce qui a changé ensuite}
    La reprise sous TwinCAT remplace cette séquence linéaire par deux niveaux : un mode global \texttt{INIT / AUTO / DEFAUT}, puis une machine d'états dédiée au cycle \texttt{AUTO}. Une PLC Visualization V1 donne accès à la simulation et à la supervision des variables exposées par \texttt{MAIN}.
  \end{notebox}

  \begin{center}
    \begin{minipage}[t]{.31\linewidth}\metric{8}{bacs mémorisés}\end{minipage}\hfill
    \begin{minipage}[t]{.31\linewidth}\metric{4}{cycles / index}\end{minipage}\hfill
    \begin{minipage}[t]{.31\linewidth}\metric{10 s}{appui long}\end{minipage}
  \end{center}
\end{minipage}

\newpage

% ==============================================================
